\cleardoubleoddpage
\begin{abstract*}
\begingroup
\enlargethispage{1\baselineskip}
\fontsize{10pt}{12pt}\selectfont
Out-of-distribution (OOD) detection is a critical challenge for safe machine learning. Deep models often face data outside their training distribution while still producing overconfident predictions. This thesis treats conditional diffusion models as generative classifiers: class-conditional reconstruction error serves as the OOD signal.

Unlike likelihood-based approaches that can be brittle in high-dimensional image spaces, diffusion models ask how well a sample can be denoised under competing class conditions. We develop a binary conditional diffusion model (CDM) for CIFAR-10 OOD detection on a within-CIFAR \emph{airplane-vs-rest} split (single ID class; multi-class generalisation is left for future work) and introduce a class-conditional separation loss that widens the reconstruction-error gap between in-distribution and out-of-distribution conditions. The central result is that the separation loss improves both performance and stability. Without the separation loss ($\lambda=0$, seeds 42, 123, and 456), the binary CDM reaches $92.52\% \pm 11.07\%$ AUROC, revealing strong seed sensitivity; with $\lambda=0.02$ over the same three seeds, performance rises to $99.03\% \pm 0.07\%$, a 6.5 percentage-point gain with far lower variance. Under the final audit policy, core claims are based on reproducible artefacts: 98.98\% AUROC on the within-CIFAR airplane-vs-rest split and 90.50--96.97\% AUROC across five externally auditable datasets (CIFAR-100, Places365, FashionMNIST, SVHN, Textures). Legacy Food-101/STL-10 values are retained for traceability only.

We further test whether the same principle transfers to industrial quality classification of inkjet print samples. We introduce a multi-head conditioning mechanism for structured manufacturing data and evaluate the public \texttt{InkjetOOD} pipeline on the public FTI\_Zer0P dataset. Under 5-fold cross-validation, the crop-based CDM reaches $0.8673 \pm 0.0230$ AUROC at $\lambda=0$. In contrast to CIFAR-10, non-zero separation-loss weights do not significantly improve performance, revealing a boundary condition: separation helps when class-conditional reconstruction distributions can be pulled apart, but it does not automatically transfer to small, fine-grained industrial datasets. Per-feature analysis shows strong heterogeneity: distance and dot features are easiest, while edge roughness features remain harder and more variable.

We analyse why diffusion-based reconstruction error succeeds where other generative approaches often fail, characterise the Monte Carlo timestep trade-off ($K=10$ provides a strong balance at about $68.6\times$ the ResNet-18 wall-clock cost on an NVIDIA GeForce RTX 2080 Ti; the $K$-AUROC curve was benchmarked at $\lambda=0.01$, and transferability to $\lambda=0.02$ was not independently verified), and provide modular public implementations using PyTorch Lightning and Hydra. Across the two evaluated settings, training objective design, specifically the separation loss, is the dominant factor in diffusion-based OOD detection, while the inkjet study shows why industrial transfer must be validated rather than assumed.

\vspace{0.08cm}
\noindent\textbf{Keywords:} Diffusion Models, Out-of-Distribution Detection, Generative Classification,\\
Industrial Quality Control, Inkjet Printing
\endgroup
\end{abstract*}
